\documentclass[11pt]{article}
\usepackage[a4paper,margin=25mm]{geometry}
\usepackage{amsmath,amssymb,amsthm,booktabs,longtable,array,xcolor}
\usepackage[colorlinks=true,linkcolor=blue!45!black,urlcolor=blue!45!black,citecolor=blue!45!black]{hyperref}
\setlength{\parskip}{5pt}
\setlength{\parindent}{0pt}
\setlength{\emergencystretch}{2em}
\theoremstyle{definition}
\newtheorem{proposition}{Proposition}[section]
\newtheorem{lemma}[proposition]{Lemma}
\newtheorem{corollary}[proposition]{Corollary}
\theoremstyle{remark}\newtheorem*{remark}{Remark}
\newcommand{\Q}{\mathbb Q}
\newcommand{\Z}{\mathbb Z}
\newcommand{\R}{\mathbb R}
\newcommand{\C}{\mathbb C}
\newcommand{\Res}{\operatorname{Res}}
\newcommand{\FP}{\operatorname{FP}}
\newcommand{\cont}{\operatorname{cont}_{\Q}}
\newcommand{\dd}{\,d}
\newcommand{\Btil}{\widetilde B}
\newcommand{\pair}{\mathcal B}
\newcommand{\rpair}{\mathcal R}
\newcommand{\ph}{\phantom{0}}
\title{One summation functional, one Gram pencil\\[3pt]
\large A common framework for the constructions in the $\zeta(5)$ manuscript}
\author{Research and exposition note}
\date{14 November 2026}
\begin{document}
\maketitle
\begin{abstract}
The most useful unification has two levels. The functionals are realizations of
one Euler--Maclaurin summation rule, with one free constant at the integer poles.
The determinant constructions are presentations of one rational bilinear form,
positive at $X=\zeta(5)$ and deformed by a residue pairing at infinity.
Local bases, prime valuations, logarithmic energy and the perturbation matrix
then become ways to estimate this same form and its determinant. We derive
these identifications, explain the precise role of Eisenstein and cotangent
kernels, and propose an explanatory organization of the proof.
\end{abstract}

\textbf{Scope.} References such as (3.1) or Appendix C refer to Aabir Fauzan's
17 September 2026 manuscript \cite{F}. The identities and elementary consequences
in this note are derived below. The manuscript's quantitative local inequalities
and its rational potential certificate are retained as inputs; this note does
not independently reverify all of those certificates. A clearer framework does
not by itself strengthen their numerical bounds.

\section{The objects to keep in the foreground}

Put $\xi=\zeta(5)$, $K=40n$, $N=3n$, $h=K-N$, and
\[
 D_m(t)=\prod_{j=1}^m(t+j^2),\qquad
 T_K(t)=\frac{D_K(t)}{D_N(t)},\qquad W_K(t)=D_N(t)^5.
\]
The two organizing objects are
\begin{equation}\label{eq:twoobjects}
 \tau_X\quad\hbox{and}\quad
 \pair_{K,X}(p,q)=\mu_X\!\left(\frac{W_Kpq}{T_K}\right),
 \qquad p,q\in E_K:=\Q[t]_{<h}.
\end{equation}
Here $\mu_X$ is a weighted pullback of $\tau_X$, not independent data:
\begin{equation}\label{eq:pullback}
 \mu_X(R)=\tau_X\bigl(x^5R(-x^2)\bigr).
\end{equation}
The matrix $G_K(X)$ is the matrix of \eqref{eq:twoobjects} in the monomial basis,
and $\Delta_K(X)=\det G_K(X)$. The equality
$D_N^6/D_K=W_K/T_K$ is cancellation in the same rational weight.

There are three useful levels of identification throughout:
\begin{itemize}
\item \emph{Exact representations}: the same functional, form or polynomial in
different coordinates.
\item \emph{Completions and estimates}: the same rational data viewed over
$\R$ or $\Q_p$, with different norms and convergence statements.
\item \emph{Auxiliary certificates}: chosen bases, allocation rules and a
comparison measure. Their roles are unified, but their choices are not forced
by one universal analytic function.
\end{itemize}

\section{Euler--Maclaurin is the source of the functional family}

\subsection{One formal operator and one boundary constant}
Let $\partial=d/dx$. The Bernoulli generating function gives the exact polynomial
operator
\begin{equation}\label{eq:EM}
 \tau(P)=\left[\frac{\partial^4}{24(e^\partial-1)}P\right](0)
 =\frac1{24}\sum_{k\ge0}\frac{B_k}{k!}P^{(k+3)}(0).
\end{equation}
The sum terminates. Thus no convergence of an infinite differential operator is
being asserted. Equivalently, as a formal generating identity,
\begin{equation}\label{eq:symbol}
 \tau(e^{ux})=\frac{u^4}{24(e^u-1)}.
\end{equation}
The manuscript's intermediate $L$, with $L(x^k)=B_k$, is simply
$[\partial/(e^\partial-1)]$ followed by evaluation at zero. Its $\tau$ is
$L\circ\partial^3/24$. This is the usual Bernoulli summation symbol
\cite{Bernoulli}, with three extra derivatives.

\begin{proposition}[An economical definition]\label{prop:axioms}
On rational functions with polynomial part and simple integer poles there is a
unique $\Q$-linear functional $\tau_X$ with values in $\Q[X]$ such that
\begin{align}
 \tau_X(g(x+1)-g(x))&=\frac{g^{(4)}(0)}{24}
                  &&\text{if $g$ is regular at zero},\label{eq:difference}\\
 \tau_X(g(-1-x))&=-\tau_X(g(x)),\label{eq:reflection}\\
 \tau_X(1/x)&=-X.\label{eq:basepole}
\end{align}
Its polynomial restriction is \eqref{eq:EM}, and
\begin{equation}\label{eq:poles}
 \tau_X\!\left(\frac1{x-r}\right)=H^{(5)}_{d(r)}-X,
 \qquad d(r)=\begin{cases}r&r\ge0,\\-r-1&r<0.\end{cases}
\end{equation}
\end{proposition}
\begin{proof}
The forward-difference map on polynomials is surjective; \eqref{eq:difference}
therefore fixes the polynomial restriction. Multiplying the symbol
\eqref{eq:symbol} by $e^u-1$ verifies it. If
$a_r=\tau_X(1/(x-r))$, the difference identity gives
$a_{r-1}-a_r=-r^{-5}$ for $r\ne0$. Starting with $a_0=-X$ gives
$a_r=H_r^{(5)}-X$ for $r\ge0$. Reflection gives $a_{-r-1}=a_r$.
Partial fractions give uniqueness and existence. The polynomial reflection
identity also follows from $e^{-u}\tau(e^{-ux})=-\tau(e^{ux})$.
\end{proof}

This explains the otherwise separate ingredients: Bernoulli moments are the
polynomial part of Euler summation, fifth harmonic sums are its increments at
simple poles, and $X$ is the one free pole-boundary value. For example,
$\tau(x^3)=1/4$, and \eqref{eq:pullback} gives
\[
 \mu(t^e)=(-1)^e\frac{(2e+3)(2e+4)(2e+5)}{24}B_{2e+2}.
\]
For a pole use
\[
 \frac{x^5}{j^2-x^2}
 =-x^3-j^2x-\frac{j^4}{2}\left(\frac1{x-j}+\frac1{x+j}\right).
\]
Equation \eqref{eq:poles} then gives exactly
$\mu_X((t+j^2)^{-1})=j^4(X-H_j^{(5)})-1/4+1/(2j)$.

\subsection{The real realization comes from the same summation kernel}
Set $f(y)=(e^{2\pi y}-1)^{-1}$. For a real rational $g$ with simple integer
poles and no pole at zero, an Abel--Plana/Hermite form of the same functional is
\begin{equation}\label{eq:AP}
 \tau_\xi(g)=\frac{g'''(0)-g^{(4)}(0)/2}{24}
 -\frac{i}{24}\int_0^\infty
 \frac{g^{(4)}(iy)-g^{(4)}(-iy)}{e^{2\pi y}-1}\dd y.
\end{equation}
This identity can be checked on polynomial monomials using Euler's formula for
even zeta values, and on the pole basis using Hermite's formula
\cite{Hurwitz,Euler}. Both sides are linear, so that check proves it on the
whole stated domain. Near zero the numerator in the integral is $O(y)$; at
infinity the exponential denominator controls every polynomial part.

For $g(x)=x^5R(-x^2)$, the boundary derivatives in \eqref{eq:AP} vanish.
Write $q(y)=y^5R(y^2)$. Then $g^{(4)}(iy)=iq^{(4)}(y)$ and
$g^{(4)}(-iy)=-iq^{(4)}(y)$. Four integrations by parts give
\begin{equation}\label{eq:positive}
 \mu_\xi(R)=\frac1{12}\int_0^\infty q(y)f^{(4)}(y)\dd y
 =\int_0^\infty R(y^2)w(y)\dd y,
\end{equation}
where
\begin{equation}\label{eq:weight}
 w(y)=\frac{y^5}{12}f^{(4)}(y)
 =\frac{(2\pi)^4y^5}{12}\sum_{\ell\ge1}\ell^4e^{-2\pi\ell y}>0.
\end{equation}
Each integration-by-parts boundary product is $O(y)$ at zero and exponentially
small at infinity. Thus positivity is a consequence of the same Euler summation
kernel, rather than an unrelated representation added later.

\subsection{One Stieltjes transform contains both definitions of \texorpdfstring{$\mu$}{mu}}
Define the renormalized Hurwitz-zeta function
\begin{equation}\label{eq:Phi}
 \Phi(a)=a^4\zeta(5,a)-\frac1{2a}-\frac14
       =\int_0^\infty\frac{w(y)}{a^2+y^2}\dd y,
 \qquad a>0.
\end{equation}
Equivalently $\zeta(5,a)=-\psi^{(4)}(a)/24$, so this is a renormalized
polygamma function \cite{Polygamma}. In the variable $z=a^2$, it is a Stieltjes
transform of the positive measure
\[
 d\nu(t)=\frac{w(\sqrt t)}{2\sqrt t}\dd t,\qquad
 \Phi(\sqrt z)=\int_0^\infty\frac{d\nu(t)}{z+t}\quad(z>0).
\]
Its integer samples and asymptotic coefficients give the two clauses of the
manuscript's definition:
\begin{align}
 \Phi(j)&=j^4(\xi-H_j^{(5)})-\frac14+\frac1{2j},\label{eq:samples}\\
 \Phi(a)&\sim\sum_{e\ge0}(-1)^e\mu(t^e)a^{-2e-2}.
 \label{eq:PhiAsymp}
\end{align}
The first is the Hurwitz recurrence. The second is the Euler--Maclaurin expansion
after its leading terms have been removed. An exact remainder, avoiding any
interpretation of a divergent expansion as an equality, is
\begin{equation}\label{eq:remainder}
 \Phi(a)=\sum_{e=0}^{m-1}\frac{(-1)^e\mu(t^e)}{a^{2e+2}}
 +\frac{(-1)^m}{a^{2m}}\int_0^\infty
       \frac{y^{2m}w(y)}{a^2+y^2}\dd y.
\end{equation}
It follows by a finite geometric identity under the integral. The last term has
sign $(-1)^m$ and absolute value at most $\mu(t^m)a^{-2m-2}$.

\subsection{A literal regularized lattice sum, if desired}
There is also one precise formula for the entire real functional, including
polynomial parts and poles at zero. If
$g=P+\sum_r c_r/(x-r)$ and $P'''(x)=\sum_k a_kx^k$, define
\[
 Z_g(z)=\sum_k a_k\zeta(-k,z)-6\sum_r c_r\zeta(4,z-r).
\]
This is exactly what is meant here by $\sum_{n\ge0}^{\zeta\text{-reg}}g'''(n+z)$:
polynomial sums are defined by Hurwitz continuation and the proper part is
summed by its convergent series, then continued. The identity
\begin{equation}\label{eq:regsum}
 \boxed{\displaystyle
 \tau_\xi(g)=-\frac1{24}\FP_{z=0} Z_g'(z)}
\end{equation}
is verified from $\zeta(-k,z)=-B_{k+1}(z)/(k+1)$ and the Hurwitz recurrence.
For a pole at $r\ge0$, the constant term of $\zeta(5,z-r)$ is
$\xi-H_r^{(5)}$. For a negative pole $r=-j$, it is $\zeta(5,j)$.

The derivative must be taken \emph{after} regularization. For $g=x^3$,
$Z_g(z)=6\zeta(0,z)=3-6z$, giving $\tau(x^3)=1/4$.
A naive sum of $g^{(4)}$ would give zero and lose this term. Formula
\eqref{eq:regsum} explains the origin, but Proposition \ref{prop:axioms} is the
lighter definition for a proof.

\section{What the Eisenstein viewpoint does explain}

\subsection{The exact kernel is one periodic lattice row}
For $\operatorname{Im}z>0$, set
\[
 \mathcal E_5(z)=\sum_{n\in\Z}(z+n)^{-5}.
\]
Differentiating the classical cotangent partial fraction four times gives
\begin{equation}\label{eq:cotrow}
 \mathcal E_5(z)=\frac1{24}\frac{d^4}{dz^4}(\pi\cot\pi z)
 =\frac{(-2\pi i)^5}{24}\sum_{\ell\ge1}\ell^4e^{2\pi i\ell z}.
\end{equation}
Consequently
\begin{equation}\label{eq:rowweight}
 \boxed{\displaystyle w(y)=\frac{i}{\pi}y^5\mathcal E_5(iy).}
\end{equation}
The Fourier expansion verifies the sign and makes positivity transparent.
Thus the positive density is a fifth-order cotangent kernel on the integer
lattice. This follows directly from the classical partial fractions
\cite{Cotangent}; it is an exact Eisenstein-type interpretation.

It is not the full weight-five level-one modular Eisenstein series.
The absolutely convergent double sum
$\sum_{(m,n)\ne(0,0)}(mz+n)^{-5}$ vanishes by pairing $(m,n)$ with $(-m,-n)$.
The single row \eqref{eq:cotrow} does not have that cancellation and is not, by
itself, a weight-five modular form.

There is a separate genuine modular connection: if $q=e^{2\pi iz}$ and
\[
 \mathcal F_5(z)=\sum_{n\ge1}\sigma_{-5}(n)q^n,
 \qquad \mathcal D=q\frac d{dq},
\]
then
\[
 \mathcal D^5\mathcal F_5=\sum_{n\ge1}\sigma_5(n)q^n
 =\frac{1-E_6(z)}{504},
\]
using the convention $E_6=1-504\sum\sigma_5(n)q^n$.
After the constant term is restored by a polynomial primitive, this is the
Eichler-integral setting in which $\zeta(5)$ occurs in a period polynomial
\cite{Ramanujan}. It explains the wider classical context. No identity found
here makes the paper's local row allocations or potential certificate a
consequence of the modular transformation law, so introducing that theory into
the main proof would add machinery without replacing those estimates.

\subsection{The products and bases also come from the integer lattice}
The same lattice supplies the finite products:
\begin{align}
 \frac{D_m(-x^2)}{(m!)^2}
 &=\prod_{j=1}^m\left(1-\frac{x^2}{j^2}\right)
 =\binom{m-x}{m}\binom{m+x}{m},\label{eq:sinprod}\\
 D_m(y^2)&=
 \frac{\Gamma(m+1+iy)\Gamma(m+1-iy)}
      {\Gamma(1+iy)\Gamma(1-iy)}.\label{eq:gammaprod}
\end{align}
For fixed $x$, \eqref{eq:sinprod} tends to $\sin(\pi x)/(\pi x)$;
at $x=iy$ the corresponding limit is $\sinh(\pi y)/(\pi y)$.
These finite sine products connect to the cotangent kernel through
\[
 \frac d{dx}\log\!\left(\frac{\sin\pi x}{\pi x}\right)
 =\pi\cot\pi x-\frac1x
\]
away from the zeros, with the removable singularity at zero understood
\cite{Cotangent,Gamma}.

The $q_i$ are normalized interpolation polynomials on the nodes
$0,-1,-4,\ldots$. Integer-valuedness follows from
\[
 q_i(-x^2)=\binom{x+i}{2i}+\binom{x+i-1}{2i}\quad(i\ge1),
 \qquad q_0=1
\]
The factors in $f_i=(D_N/(N!)^2)^3q_i$ absorb half
of the sixth-power weight into each argument of the bilinear form.
Stirling/Euler summation and Legendre's factorial-valuation formula later
measure this same product data at the real and finite places.

\section{One Gram pencil and one residue pairing}

\subsection{The parameter changes only a Laurent coefficient}
For rational $R$ in the manuscript's domain, write $[t^{-3}]_\infty R$ for its
Laurent coefficient at infinity. Partial fractions immediately give
\begin{equation}\label{eq:mastermu}
 \boxed{\displaystyle
 \mu_X(R)=\int_0^\infty R(y^2)w(y)\dd y
             +(X-\xi)[t^{-3}]_\infty R.}
\end{equation}
Indeed the coefficient of $t^{-3}$ in $(t+j^2)^{-1}$ is $j^4$, and polynomial
parts contribute zero. Equivalently this last functional is
$-\Res_{t=\infty}t^2R(t)\dd t$.
Similarly $\tau_X(g)=\tau_\xi(g)-(X-\xi)[x^{-1}]_\infty g$.

Define
\[
 d\nu_K(t)=\frac{W_K(t)}{T_K(t)}d\nu(t),\qquad
 \rpair_K(p,q)=[t^{-1}]_\infty\frac{t^2W_K(t)p(t)q(t)}{T_K(t)}.
\]
Then the promised common object is
\begin{equation}\label{eq:masterform}
 \boxed{\displaystyle
 \pair_{K,X}(p,q)=\int_0^\infty p(t)q(t)\,d\nu_K(t)
                     +(X-\xi)\rpair_K(p,q).}
\end{equation}
The first form is positive definite on $E_K\otimes\R$. The second is rational
and nondegenerate. Although the displayed real splitting contains $\xi$, the
original $\pair_{K,X}$ has coefficients in $\Q[X]$ by the functional formulas.
This is a one-parameter deformation, not a rank-one perturbation: the residue
matrix has rank $h$, as the next calculation proves.

\subsection{The slope is a finite residue algebra}
Let $\mathcal A_K=\Q[t]/(T_K)$, and let $\epsilon$ extract the coefficient of
$t^{h-1}$ in the unique remainder of degree less than $h$. Put $H_K=t^2W_K$.
Then
\begin{equation}\label{eq:frobenius}
 \rpair_K(p,q)=\epsilon(H_Kpq)
 =\sum_{j=N+1}^K c_jp(-j^2)q(-j^2),\quad
 c_j=\frac{j^4D_N(-j^2)^5}{T_K'(-j^2)}.
\end{equation}
The multiplier $H_K$ is a unit in $\mathcal A_K$ because it has no common root
with $T_K$. The matrix of $\epsilon(pq)$ has zero entries when $i+j<h-1$
and ones when $i+j=h-1$, with indices starting at zero. Its determinant is
$(-1)^{h(h-1)/2}$. Multiplication by $H_K$ has determinant
$\Res(T_K,H_K)$. Therefore
\begin{equation}\label{eq:resultant}
 \det\rpair_K=(-1)^{h(h-1)/2}\Res(T_K,H_K)
 =(-1)^{h(h-1)/2}\prod_{j=N+1}^K j^4D_N(-j^2)^5\ne0.
\end{equation}
This proves the degree-$h$ assertion and reproduces the manuscript's (2.9).
The same residue form is $M_K=[X]G_K(X)$ in Appendix C.

\subsection{One determinant identity explains degree and positivity}
Equation \eqref{eq:masterform} is the moment form of the signed measure
\[
 d\nu_{K,X}=d\nu_K+(X-\xi)\sum_{j=N+1}^K c_j\delta_{-j^2}.
\]
Andr\'eief's identity \cite{Andreief} gives the polynomial identity
\begin{equation}\label{eq:andreief}
 \Delta_K(X)=\frac1{h!}\int_{\R^h}
       \prod_{i<j}(t_i-t_j)^2\prod_{i=1}^h d\nu_{K,X}(t_i).
\end{equation}
For real $X$ this is an ordinary finite signed-measure identity; expansion
gives its polynomial interpretation. At $X=\xi$ all measures are positive,
giving strict positivity. The leading coefficient selects all $h$ atomic
factors. Repeated nodes give zero; the $h!$ permutations of the $h$ distinct
nodes give exactly \eqref{eq:resultant}. Thus the real nonvanishing argument and
the degree argument are two parts of the same identity.

\subsection{Appendix C is the spectral theory of the same pencil}
Let $G=G_K(\xi)$, $M=M_K$ and $B=G^{-1/2}MG^{-1/2}$, with $\|B\|$
the Euclidean operator norm. Then
\begin{equation}\label{eq:pencil}
 G_K(\xi+\delta)=G^{1/2}(I+\delta B)G^{1/2},\qquad
 \frac{\Delta_K(\xi+\delta)}{\Delta_K(\xi)}
     =\prod_{j=1}^h(1+\delta\beta_j),
\end{equation}
where $\beta_j$ are the real eigenvalues of $B$. Nondegeneracy of $M$ implies
that none is zero.
\begin{corollary}\label{cor:roots}
All roots of $\Delta_K$ are real, with multiplicities allowed, and
\[
 \operatorname{dist}(\xi,\{X:\Delta_K(X)=0\})=\|B\|^{-1}.
\]
\end{corollary}
This is a consequence of the manuscript's pencil, not an additional assumption.
It makes the purpose of the relative norm estimate transparent: it is a
root-exclusion estimate. Equation \eqref{eq:pencil} also bounds the determinant
near $\xi$ and the coefficient height, so those are two consequences of one
spectral estimate.

Orthogonal polynomials give an optional exact description. If $p_0,\ldots,p_{h-1}$
are orthonormal for $\nu_K$, then the relative residue matrix is
$\sum_j c_j v_jv_j^T$, where $v_j=(p_i(-j^2))_i$. Hence
\[
 \|B\|\le\sum_j|c_j|\mathcal K_h(-j^2,-j^2),\qquad
 \mathcal K_h(s,s)=\sum_{i=0}^{h-1}p_i(s)^2.
\]
Appendix C bounds this kind of polynomial evaluation relative to its integral
norm using Legendre polynomials. Introducing the full orthogonal-polynomial
theory is optional; the Gram-pencil identity already supplies the needed logic.

\section{Local arithmetic from Volkenborn summation}

\subsection{The \texorpdfstring{$p$}{p}-adic realization of the Euler functional}
For polynomials define the Volkenborn integral
\[
 I_p(P)=\lim_{m\to\infty}p^{-m}\sum_{a=0}^{p^m-1}P(a),
\]
where the limit is in $\Q_p$. Its moments are $I_p(x^k)=B_k$, so
\begin{equation}\label{eq:Volkenborn}
 \tau(P)=\frac1{24}I_p(P''').
\end{equation}
For polynomials $f$, the translation defect is
$I_p(f(x+1))-I_p(f(x))=f'(0)$.
Splitting into residue classes and accounting for three derivatives gives
\[
 \tau(P)=p^{-4}\sum_{a=0}^{p-1}\tau(P(a+px)).
\]
These are standard Volkenborn identities \cite{Volkenborn,HuKim}; the rational
value in \eqref{eq:Volkenborn} is independent of $p$.
This integral is not translation-invariant Haar integration, and rational
functions with poles in $\Z_p$ are not covered by its ordinary definition.
The manuscript uses the algebraic extension of Proposition \ref{prop:axioms}.

For $p\ge7$, the moments of $\tau$ have valuations at least $-1$, and are
integral through degree $p+1$. This permits continuous evaluation on the Tate
algebra. A far pole is represented by the convergent expansion
\begin{equation}\label{eq:farpole}
 \tau^{\rm an}\!\left(\frac1{x+a}\right)
 =-\frac14\sum_{k\ge0}(-1)^k\binom{k+3}{3}B_k a^{-k-4},
 \qquad |a|_p>1.
\end{equation}
The very same formal Laurent series is the real Euler--Maclaurin asymptotic
expansion of $-\zeta(5,a)$. Real asymptotic summation and $p$-adic convergence
are different realizations of the same coefficients, not equal numerical
values in one field.

For near poles, append their principal parts to the analytic space and use
$\tau_Y(1/(x-r))=H_{d(r)}^{(5)}-Y$. This is $\tau_Y^{\rm ext}$.
Demanding distribution at $1/x$ forces
\begin{equation}\label{eq:param}
 C_p=\sum_{a=1}^{p-1}\tau^{\rm an}\!\left(\frac1{x+a/p}\right),\quad
 Y=p^5X+C_p,\quad
 \tau_X(g)=p^{-4}\sum_{a=0}^{p-1}\tau_Y^{\rm ext}(g(a+px)).
\end{equation}
Indeed the right side at $1/x$ is $p^{-5}(-Y+C_p)=-X$.
The fifth power is forced by four powers from the functional and one from
rescaling the pole. Reflection and the difference law propagate the identity
to the other poles. The low-degree integrality condition and the estimate
$C_p\in p^5\Z_p$ remain substantive parts of the local proof. Lemma 3.1 also
requires unit-separated near poles and harmonic indices $d(r)<p$ so that the
principal-part values and partial-fraction decomposition are integral.

\subsection{The three bases estimate the same pairing}
Every change of basis $C$ satisfies
\begin{equation}\label{eq:congruence}
 G_{\rm new}=C^TG_{\rm old}C,\qquad
 \det G_{\rm new}=(\det C)^2\det G_{\rm old}.
\end{equation}
The choices in the manuscript have distinct practical advantages:
\begin{center}
\begin{tabular}{>{\raggedright\arraybackslash}p{.19\textwidth}>{\raggedright\arraybackslash}p{.28\textwidth}>{\raggedright\arraybackslash}p{.43\textwidth}}
\toprule
Prime range & Basis description & Purpose\\ \midrule
Small & Integer-valued $q_i$; weighted $f_i$ & Use finite differences and the
Bernoulli functional with coarse uniform losses.\\[4pt]
Inner & $E_{a,i}$, with prescribed jets at $t=-a^2$ & Allocate vanishing so that
every residue-class contribution respects assigned row weights.\\[4pt]
Outer & $P_aq_{a,i}$ using exact class factors of $T_K$ & Cancel dangerous poles,
separate residue blocks, and retain a controlled polynomial correction.\\
\bottomrule
\end{tabular}
\end{center}
The two local bases are unimodular. The small-prime calculation also rescales
the pairing and incorporates its explicit determinant factor $S_K$.
Precisely, the $q_i$ form a basis of $E_K$, whereas the $f_i$ are their images
in the weighted space $D_N^3E_K$ under multiplication by
$D_N^3/(N!)^6$. They are not literally a basis of $E_K$.
The outer construction is the Chinese-remainder decomposition of the actual
pole algebra $\Z_p[t]/(T_K)$. The inner construction instead uses an auxiliary
jet algebra with modulus $\prod_a(t+a^2)^{L_a}$. They are bases of the same
degree-$<h$ polynomial space, but these two quotient algebras need not be equal.

The common determinant mechanism is elementary. If
$v_p^G(A_{ij})\ge w_i+w_j$, every determinant term has valuation at least
$2\sum_iw_i$. In the outer range the full form is
$A+p^{-1}L$ with $L$ integral and rank at most $r$. When the $w_i$ are
nonpositive half-integers and exactly $z$ are zero, complementary minors give
\begin{equation}\label{eq:rank}
 v_p^G\det(A+p^{-1}L)\ge2\sum_iw_i-\min(r,z).
\end{equation}
Minors of $L$ above rank $r$ vanish. If the weights are sorted in decreasing
order, the possible loss is bounded by
$\max_{0\le k\le r}(k+2\sum_{i=1}^k w_{(i)})$, with $k\le h$ understood.
Its increments are one at zero weights and nonpositive thereafter.
This gives $\min(r,z)$ and explains the correction.

The finite residue algebra describes the slope and residue terms. It does not
erase polynomial quotients in the full functional. Their contribution is
precisely why the rank correction and the analytic completion cannot be dropped.

\subsection{The profiles are sums of the same local weights}
The floor-function count is an exact rescaling:
\[
 \ell_A(a)=\ell(A/p,a/p),\qquad
 \ell(x,z)=\lfloor x-z\rfloor+\lfloor x+z\rfloor+1.
\]
The later functions summarize earlier local inequalities:
\begin{align*}
 \gamma_p^{\rm in}&=p\Gamma(K/p)+O_M(1),&
 v_p(S_K)&=p\mathcal N(K/p)+O_M(1),\\
 -\gamma_p^{\rm out}&=K(R_0(p/K)-d(p/K))+O(1),&
 R(x)&=-\Gamma(x)-\mathcal N(x).
\end{align*}
Here $\mathcal N$ denotes the manuscript's function $N(x)$, to distinguish it
from the integer $N$. The prime number theorem turns these valuation sums
into the integrals of Section 5. These functions are accounting profiles,
not new moment functionals.

There is a second exact Euler--Maclaurin appearance in the tail estimate.
Let $\alpha=3/40$, $\lambda=37/40$, and
$\Btil_k(x)=B_k(\{x\})$. The manuscript's $F,P,C$ satisfy
\begin{align}
 F(x)+\lambda&=2\lambda\Btil_1(x)-12\lambda\Btil_1(\alpha x),\label{eq:tail1}\\
 P(x)-\overline P&=\lambda\Btil_2(x)-74\Btil_2(\alpha x),\label{eq:tail2}\\
 C(x)&=\frac\lambda3\Btil_3(x)-\frac{74}{3\alpha}\Btil_3(\alpha x),
 \qquad \overline P=\frac{2923}{240}.\label{eq:tail3}
\end{align}
Since $74=6\lambda/\alpha$, differentiation gives $P'=F+\lambda$ and
$C'=P-\overline P$ away from the breakpoints. The periodic functions $P,C$
are thus Bernoulli primitives of the lattice-counting error. The two
integrations by parts in (5.15) are the usual Euler--Maclaurin treatment of
oscillatory counting errors. This is the same Bernoulli machinery as
\eqref{eq:EM}, acting on a different task.

\section{At the real place: determinant energy and its certificate}

At $X=\xi$, \eqref{eq:andreief} is a positive Vandermonde integral. Under
$y_i=K\sqrt{t_i}$, its configuration-dependent leading logarithm, after
extracting deterministic powers of $K$, is
\begin{equation}\label{eq:discreteenergy}
 2\sum_{i<j}\log|t_i-t_j|-K\sum_iV(t_i),
\end{equation}
where
\[
 V(t)=2\pi\sqrt t+\int_0^1\log(t+u^2)\dd u
                   -6\int_0^\alpha\log(t+u^2)\dd u.
\]
The field is the continuum logarithm of the original rational weight, including
the exponential decay of $w$. The sum-to-integral step is uniform even at $t=0$:
if
\[
 E_m(t)=\sum_{j=1}^m\log(t+(j/K)^2)
       -K\int_0^{m/K}\log(t+u^2)\dd u,
\]
then $E_m(t)\ge0$ by right Riemann sums, $E_m'(t)\le0$ for $t>0$, and
\[
 E_m(0)=2\log(m!)-2m\log m+2m=O(1+\log m).
\]
This justifies the field directly, without a nonuniform use of Stirling near an
endpoint. The determinant and its energy are the same object at different
scales, with the error terms kept in the proof.

For a positive measure $\eta$ of mass $\lambda=h/K$, write
$U^\rho(t)=\int\log|t-u|\,d\rho(u)$,
$I(\eta)=\iint\log|t-u|\,d\eta(t)d\eta(u)$ and
$\mathcal E_V(\eta)=I(\eta)-\int V\,d\eta$.
The key certificate lemma is short.
\begin{lemma}[A dual upper bound]\label{lem:energy}
Suppose $\eta,\rho$ are positive measures of mass $\lambda$, with finite energies
in the expressions below, and
$2U^\rho(t)-V(t)\le M_0$ on the allowed support of $\eta$.
If $I(\eta-\rho)\le0$, then
\[
 \mathcal E_V(\eta)\le\lambda M_0-I(\rho).
\]
\end{lemma}
\begin{proof}
Expand
$I(\eta-\rho)=I(\eta)-2\int U^\rho\,d\eta+I(\rho)$ and use the potential
inequality. The equal masses are what permit the zero-mass energy inequality.
\end{proof}
For point configurations the diagonal energy is infinite, so the paper replaces
each point by a small circle and uses
$\sigma=K^{-1}\sum_i\omega_i$. This is a regularization of the configuration,
not a replacement of the original moment measure. The sixteen-arcsine $\rho$
is a chosen \emph{dual certificate}: it need not be the exact equilibrium
measure, and it is not $\nu_K$. The circle error and the retained integrable
tail still have to be controlled as in Section 6. In detail, the zero-mass
energy theorem applies to a compactly supported real signed measure $\theta$
on $\C$ with
\[
 \theta(\C)=0,\qquad
 \iint\bigl|\log|z-w|\bigr|\,d|\theta|(z)d|\theta|(w)<\infty.
\]
The real-axis inequality $2U^\rho-V\le M_0$ cannot simply be applied on the
complex circle support of $\sigma$. The source instead proves
\[
 \int U^\rho\,d\omega_i\le U^\rho(t_i)+60\sqrt\varepsilon,
 \qquad I(\omega_i)=\log\varepsilon,
\]
and bounds mutual circle energies below by $\log|t_i-t_j|$. Applied at the
real centers, these give the necessary error terms
\begin{align*}
 2\sum_{i<j}\log|t_i-t_j|-K\sum_iV(t_i)
 &\le (\lambda M_0-I(\rho))K^2\\
 &\quad+120\lambda K^2\sqrt\varepsilon-h\log\varepsilon.
\end{align*}
The abstract certificate lemma does not remove this regularization step.
This is standard logarithmic
potential theory in the particular setting of the manuscript \cite{Saff}.

\section{One invariant for all the normalizations}

For a nonzero $A\in\Q[X]$ with $A(\xi)\ne0$, define its primitive evaluation size
\begin{equation}\label{eq:J}
 \mathcal J_\xi(A)=\log|A(\xi)|-
                  \sum_p v_p^G(A)\log p
 =\log\left|\frac{A(\xi)}{\cont(A)}\right|.
\end{equation}
The sum is finite. The positive rational content has
$v_p(\cont(A))=v_p^G(A)$, and $A/\cont(A)$ is a primitive integer polynomial.
For every nonzero rational $c$,
\begin{equation}\label{eq:invariant}
 \mathcal J_\xi(cA)=\mathcal J_\xi(A),
\end{equation}
since $\log|c|=\sum_pv_p(c)\log p$. Thus $\Delta_K,F_K,Q_{K,M},P_K$
all have the same primitive evaluation size. They are different explicit
scalings of one polynomial line.

This is an elementary form of local-to-global bookkeeping. Its finite terms
are valuations of rational coefficients; it does not assign $p$-adic values
to the real number $\xi$, or apply a number-field product formula to it.
It is also an evaluation size, not the usual coefficient height.

Here $F_K=S_K\Delta_K$, $Q_{K,M}=m_{K,M}F_K$, and
$P_K=\Delta_K/\cont(\Delta_K)$, using the paper's notation.
The factorial scalar $S_K$ is useful because it cancels the $K^2\log K$ term
in the real estimate and is adapted to the small-prime basis. It cancels
\emph{exactly} from \eqref{eq:J}. The remaining scalar
$m_{K,M}=\prod_p p^{-L_p}$ is a finite product using certified bounds
$L_p\le v_p^G(F_K)$; take $L_p=0$ outside it only where integrality is proved.
For $Q=m_{K,M}F_K\in\Z[X]$ and $F_K(\xi)>0$,
\begin{equation}\label{eq:certified}
 \log Q(\xi)-\mathcal J_\xi(F_K)
 =\sum_p\bigl(v_p^G(F_K)-L_p\bigr)\log p
 =\log\cont(Q)\ge0.
\end{equation}
So the explicit normalization is an upper-bound certificate for the canonical
primitive size. Positive $L_p$ allow division by primes; negative ones clear
denominators. Discarding their signs would discard genuine savings.

The abstract finishing principle can now be stated without any of the
auxiliary names.
\begin{proposition}[Primitive evaluation criterion]\label{prop:criterion}
Suppose $A_K\in\Q[X]$ has degree $O(K)$, $A_K(\xi)\ne0$, and
$\mathcal J_\xi(A_K)\le-cK^2+o(K^2)$ for some $c>0$.
Then $\xi$ is irrational.
\end{proposition}
\begin{proof}
If $\xi=a/b$ with $b>0$, and $d_K=\deg A_K$, then
$b^{d_K}A_K(a/b)/\cont(A_K)$ is a nonzero integer. Consequently
$\mathcal J_\xi(A_K)\ge-d_K\log b$, contradicting the quadratic bound.
\end{proof}
In the manuscript the estimates are
$\log F_K(\xi)\le UK^2+o(K^2)$ and
$\log m_{K,M}\le A_MK^2+o(K^2)$, with $A_M+U<0$ at the chosen cutoff.
The irrationality proof and the approximation proof use the same integer
criterion; the latter first uses \eqref{eq:pencil} to retain nonvanishing at
nearby rational points.

\section{A dictionary and a proposed presentation}

\begin{longtable}{>{\raggedright\arraybackslash}p{.27\textwidth}>{\raggedright\arraybackslash}p{.66\textwidth}}
\toprule
Manuscript construction & Common object or operation\\ \midrule
\endfirsthead
\toprule Manuscript construction & Common object or operation\\ \midrule
\endhead
$L,\tau,\tau_X$ & Euler--Maclaurin operator; differentiated Bernoulli moments;
one-parameter singular extension.\\[4pt]
$\mu_X$ & Weighted square-coordinate pullback of $\tau_X$.\\[4pt]
Pole values involving $H_j^{(5)}$; Bernoulli moments & Integer samples and asymptotic coefficients
of the same Hurwitz/Stieltjes function $\Phi$.\\[4pt]
$w$ & Fourth derivative of the Euler summation kernel; equivalently a
fifth-order cotangent row.\\[4pt]
$\tau^{\rm an},\tau_Y^{\rm ext}$ & Continuous $p$-adic evaluation of the same
Bernoulli data, plus specified principal parts.\\[4pt]
$Y,C_p$ & Forced change of the single boundary constant under residue-class
dilation.\\[4pt]
$D_m,W_K,T_K$ & Finite integer-lattice products and cancellation of one common
rational weight.\\[4pt]
$G_K,\Delta_K$ & Matrix and determinant of one bilinear pencil on $E_K$.\\[4pt]
$[X]G_K,M_K$ & Residue-at-infinity pairing on $\Q[t]/(T_K)$; its determinant is
a signed resultant.\\[4pt]
$q_i,f_i,E_{a,i},P_aq_{a,i}$ & Frames adapted to integer-valuedness, jets or exact
pole cancellation; all estimate the same bilinear form.\\[4pt]
$w_{a,i},\gamma_p^{\rm in/out}$ & Entrywise valuation budgets and their induced
determinant bounds, including a rank correction.\\[4pt]
$\ell,\Gamma,\mathcal N,R,R_0,d$ & Rescaled counts and limiting profiles of those
same local budgets.\\[4pt]
Tail $F,P,C$ & Periodic Bernoulli errors and their successive primitives.\\[4pt]
Vandermonde integral, $V,I$ & Determinant identity and its scaled logarithmic
energy.\\[4pt]
$\sigma,\rho$ & Regularized empirical configuration and an independently chosen
equal-mass dual certificate.\\[4pt]
$S_K,m_{K,M},F_K,Q_{K,M},P_K$ & Normalizing scalars and resulting representatives
of one polynomial line, compared by $\mathcal J_\xi$.\\[4pt]
$B_K$, eigenvalue product & The residue form in coordinates orthonormal for the
positive form; root exclusion and coefficient control.\\
\bottomrule
\end{longtable}

For a revised exposition, the following sequence would make the dependencies
visible before introducing the technical formulas.
\begin{enumerate}
\item \textbf{State the target criterion.} Prove Proposition
\ref{prop:criterion} first: the task is a nonzero determinant of degree $O(K)$
whose primitive evaluation size decays quadratically.
\item \textbf{Introduce one Euler functional.} Use Proposition
\ref{prop:axioms}; derive $\mu_X$ by weighted pullback. Present \eqref{eq:Phi}
as the explanation for the rational moments and pole values.
\item \textbf{Prove the master Gram-pencil lemma.} Establish
\eqref{eq:masterform}--\eqref{eq:resultant} and \eqref{eq:andreief} together.
Degree, positivity, the residue matrix, and the later perturbation argument
then have a single source.
\item \textbf{Prove one finite-place bound.} Explain the three ranges as three
adapted frames for this form. Isolate analytic completion and the
weight-plus-rank determinant lemma before doing the class allocations.
Put long allocation tables and exact prime-integral arithmetic in appendices.
\item \textbf{Prove one real-place bound.} Derive $V$ from the rational weight;
state Lemma \ref{lem:energy}; then supply the explicit $\rho$ and its interval
certificate. State the regularization and tail errors explicitly.
\item \textbf{Compare two costs.} Use \eqref{eq:certified}, the real estimate
and the prime sum to obtain the negative quadratic coefficient.
\item \textbf{Make Appendix C a Gram-pencil corollary.} Introduce the relative
norm as a distance-to-roots bound, with approximation and coefficient height
as its two applications.
\end{enumerate}

The cotangent/Eisenstein discussion and the regularized half-lattice formula
are useful motivation or appendices. They need not be prerequisites for the
main proof. Likewise, orthogonal polynomials give good interpretations but do
not automatically supply the necessary arithmetic estimates.

The explanatory gain is that one summation rule supplies the functional
family, and one Gram pencil supplies the determinant family. This is a
reorganization of structures largely already present in the manuscript,
not a claim that a new analytic identity replaces its quantitative work.
The local inequalities, their degree restrictions, their rank correction,
uniform limiting errors, and the final sign comparison remain essential.

\subsection*{Assessment after independent audit}
Fresh reviews rederived the functional, determinant and arithmetic
identifications. They found no substantive error in the central identities.
The normalized sine-product logarithmic derivative, the domain of the
Volkenborn translation law, and the energy-regularization hypotheses have
been made explicit in this revision.

The framework is best judged by explanation and checkability, rather than
length. Its strongest benefits are that the functional axioms force the
harmonic increments, the exact Hurwitz integral connects the two kinds of
moment data, the residue slope makes Appendix C natural, and the primitive
evaluation size exposes the signs in normalization. The asymptotic Bernoulli
series alone would not determine $\Phi$ or establish positivity: the exact
Hermite identity is indispensable. The regularized half-lattice formula is
exact under the stated prescription, but does not supply a new uniqueness
principle for regularization.

Several layers should remain distinct. The order of the summation functional
encodes the target $\zeta(5)$. The numerator power six, the ratio $N/K=3/40$,
the local bases and the arcsine comparison measure are design choices for
balancing the two estimates; they are not forced by the Euler kernel.
The finite residue algebra controls the slope and residue terms, not the
entire positive form. The comparison measure is a certificate, not the
moment measure or an asserted equilibrium measure. These distinctions make
the argument more robust against unjustified generalizations.

One optional logical economy is worth recording. The main irrationality
argument only needs $\deg\Delta_K\le h$, immediate from affine entries, and
$\Delta_K(\xi)>0$, supplied by the Gram form. Thus the exact-degree/resultant
calculation can be a structural corollary rather than a prerequisite for
irrationality. The same degree upper bound suffices for Appendix C's integer
and height arguments. If the stronger exact-degree statement is retained,
its proof remains necessary. Keeping the residue calculation for its
explanatory value is entirely consistent with this dependency distinction.

\section*{Checks and reproducibility}
The companion standard-library script \texttt{framework\_checks.py} verifies
the polynomial translation, reflection and distribution identities through
degree 24; the weighted moment identity through degree 11; and the residue
pairing/resultant formula in three exact finite examples. It checks the
periodic Bernoulli formulas at rational points and compares the Hurwitz and
Stieltjes expressions, and the cotangent-row density, numerically. The output
is \texttt{framework\_checks.json}. These are checks on the new identifications,
not substitutes for the proofs above or for the original manuscript's global
certificates. Separate research notes retain the three subagents' derivations
for the functional, arithmetic and determinant families.

\begin{thebibliography}{99}
\bibitem{F} A. Fauzan, \emph{$\zeta(5)$ is irrational}, manuscript dated
17 September 2026, v1. \url{https://zenodo.org/records/22826419}.
The local copy is \texttt{sources/Fauzan\_2026\_zeta5.pdf}.
\bibitem{Bernoulli} NIST DLMF, Bernoulli generating functions and identities,
\url{https://dlmf.nist.gov/24.2} and \url{https://dlmf.nist.gov/24.4}.
\bibitem{Hurwitz} NIST DLMF, Hurwitz zeta, recurrence, Euler--Maclaurin and
Hermite representations, \url{https://dlmf.nist.gov/25.11}.
\bibitem{Euler} NIST DLMF, even zeta values,
\url{https://dlmf.nist.gov/25.6.E2}.
\bibitem{Polygamma} NIST DLMF, polygamma functions,
\url{https://dlmf.nist.gov/5.15}.
\bibitem{Cotangent} NIST DLMF, sine product and cotangent partial fractions,
\url{https://dlmf.nist.gov/4.22}.
\bibitem{Gamma} NIST DLMF, gamma recurrence and reflection,
\url{https://dlmf.nist.gov/5.5}.
\bibitem{Ramanujan} B. C. Berndt and A. Straub, \emph{Ramanujan's Formula
for $\zeta(2n+1)$}, especially Section 5,
\url{https://arminstraub.com/downloads/pub/ramanujanzeta.pdf}.
\bibitem{Volkenborn} A. Volkenborn, \emph{On generalized $p$-adic integration},
M\'emoires de la Soci\'et\'e Math\'ematique de France 39--40 (1974), 375--384,
\url{https://www.numdam.org/book-part/MSMF_1974__39-40__375_0/}.
\bibitem{HuKim} S. Hu, D. Kim and M.-S. Kim, \emph{New identities involving
Bernoulli, Euler and Genocchi numbers}, Advances in Difference Equations 2013,
74, Section 2, \url{https://doi.org/10.1186/1687-1847-2013-74}.
\bibitem{Andreief} P. J. Forrester, \emph{Meet Andr\'eief, Bordeaux 1886,
and Andreev, Kharkov 1882--83},
\url{https://arxiv.org/abs/1806.10411}.
\bibitem{Saff} E. B. Saff, \emph{Logarithmic Potential Theory with Applications
to Approximation Theory}, \url{https://www.math.vanderbilt.edu/saffeb/texts/230.pdf}.
\end{thebibliography}
\end{document}
